\chapter{코드 전체 목록}
\label{app:code}

이 부록은 1권에서 구현한 핵심 모듈의 전체 코드를 모은 것이다. 본문에서는 설명을 위해 부분적으로 발췌했지만, 여기서는 완전한 코드를 제공한다. 최신 버전은 GitHub 저장소에서 확인하라.

\section{프로젝트 구조}

\begin{lstlisting}[style=bashstyle]
make-llm-basic/
  src/makellm/
    __init__.py
    tokenizer/
      __init__.py
      base.py        # Tokenizer 추상 클래스
      bpe.py         # BPE 토크나이저
      unigram.py     # Unigram 토크나이저
    model/
      __init__.py
      embedding.py   # TokenEmbedding, PositionalEmbedding
      attention.py   # MultiHeadAttention, CausalMask
      transformer_block.py
      gpt.py         # GPT 모델
    training/
      __init__.py
      dataset.py     # LMDataset
      loss.py        # cross_entropy_loss
      optimizers.py  # AdamW + cosine scheduler
      trainer.py     # Trainer 클래스
    inference/
      __init__.py
      sampler.py     # Greedy, Top-K, Top-P, Temperature
      generate.py    # generate 함수
    utils/
      __init__.py
      config.py      # ModelConfig, TrainerConfig
      seed.py        # set_seed
      logging.py     # Logger
  tests/             # pytest 단위 테스트
  book/              # 이 책의 LaTeX 소스
\end{lstlisting}

\section{테스트 실행}

\begin{lstlisting}[style=bashstyle]
# 패키지 설치
cd make-llm-basic
pip install -e .

# 전체 테스트 실행
pytest tests/ -v

# 특정 모듈만 테스트
pytest tests/test_tokenizer.py -v
pytest tests/test_model.py -v
pytest tests/test_training.py -v
\end{lstlisting}

\section{빠른 시작 예제}

\begin{lstlisting}[language=Python]
"""1권의 모든 모듈을 사용한 end-to-end 예제."""
from makellm.tokenizer import BPETokenizer
from makellm.model import GPT
from makellm.training import LMDataset, Trainer
from makellm.inference import generate, TopPSampler
from makellm.utils import ModelConfig, TrainerConfig, set_seed

set_seed(42)

# 1. 토크나이저 학습
corpus = ["the quick brown fox jumps", "the lazy dog sleeps"] * 100
tokenizer = BPETokenizer()
tokenizer.train(corpus, vocab_size=500)

# 2. 모델 생성
model_config = ModelConfig(
    vocab_size=tokenizer.vocab_size,
    context_length=32, d_model=64, n_heads=4, n_layers=2, d_ff=128
)
model = GPT(model_config)
print(f"Model: {model}")
print(f"Parameters: {model.num_parameters():,}")

# 3. 데이터셋 준비
dataset = LMDataset(corpus, tokenizer, context_length=32)

# 4. 학습
trainer_config = TrainerConfig(
    batch_size=4, learning_rate=3e-4, max_steps=100,
    warmup_steps=10, device="cpu"
)
trainer = Trainer(model, dataset, trainer_config, model_config)
trainer.train()

# 5. 텍스트 생성
sampler = TopPSampler(p=0.9, temperature=0.8)
text = generate(model, tokenizer, prompt="the", max_new_tokens=20, sampler=sampler)
print(f"Generated: {text}")
\end{lstlisting}

\section{모듈별 API 요약}

\begin{center}
\begin{tabularx}{\textwidth}{lX}
\toprule
\textbf{모듈} & \textbf{주요 API} \\
\midrule
\code{makellm.tokenizer} & \code{BPETokenizer, UnigramTokenizer} \\
\code{makellm.model} & \code{GPT, MultiHeadAttention, TransformerBlock} \\
\code{makellm.training} & \code{LMDataset, Trainer, cross\_entropy\_loss} \\
\code{makellm.inference} & \code{generate, GreedySampler, TopKSampler, TopPSampler} \\
\code{makellm.utils} & \code{ModelConfig, TrainerConfig, set\_seed} \\
\bottomrule
\end{tabularx}
\end{center}

\section{설정 데이터클래스}

\begin{lstlisting}[language=Python]
@dataclass
class ModelConfig:
    vocab_size: int = 1000
    context_length: int = 128
    d_model: int = 128
    n_heads: int = 4
    n_layers: int = 4
    d_ff: int = 512
    dropout: float = 0.1
    use_bias: bool = True
    layer_norm_eps: float = 1e-5
    pad_token_id: int = 0

@dataclass
class TrainerConfig:
    batch_size: int = 32
    learning_rate: float = 3e-4
    weight_decay: float = 0.1
    max_epochs: int = 10
    warmup_steps: int = 100
    lr_scheduler: str = "cosine"
    grad_clip: float = 1.0
    log_every: int = 10
    eval_every: int = 200
    save_every: int = 1000
    seed: int = 42
    device: str = "cpu"
    out_dir: str = "./runs"
\end{lstlisting}

이 설정들을 조정하여 다양한 크기의 모델을 실험할 수 있다. 예를 들어 GPT-2 small과 비슷한 설정은 \code{d\_model=768, n\_heads=12, n\_layers=12, context\_length=1024}다 (다만 실제 학습에는 상당한 컴퓨팅 자원이 필요하다).

\section{라이선스}

코드는 MIT 라이선스로 배포된다. 자유롭게 복사, 수정, 상업적 사용이 가능하다. 책 본문은 CC BY-NC-SA 4.0 (저자 확인 필요).
